% Part: proof-theory
% Chapter: proof-search
% Section: introduction

\documentclass[../../../include/open-logic-section]{subfiles}

\begin{document}

\olfileid{pt}{ps}{int}

\olsection{引言}

像 \Log{G3c} 或 \Log{N1i} 这样的推导系统，其公理和推理规则定义了哪些相继式树或公式树算作正确的推导。给定这样一棵相继式或公式树（可能还附带诸如每步应用了哪条规则、或为假设以及解除假设的推理规则所加的标记等额外信息），公理和规则的定义使我们能检查给定的树是否为正确的推导。而\emph{寻找}一个给定相继式的推导，或者从给定假设出发寻找给定公式的推导，则是一个不同的、通常也更困难的问题。这就是\emph{证明搜索}问题。

有一个非常简单粗暴的证明搜索过程。我们可以将公式看作来自有限字母表的符号序列。相继式以及相继式或公式的树也可以视为公式的序列（或许要用特殊的括号来表示树的分支）。公理和规则使我们能够机械地检查一个符号序列是否表示一个正确的推导。因此，我们可以机械地生成\emph{所有可能的}正确推导（通过生成所有用该字母表表示推导的符号序列，并检查每个候选序列是否代表一个正确推导）。这种简单粗暴的证明搜索过程就是：生成所有正确的推导，直到我们找到一个推导，它是给定相继式（或者是其未解除假设包含在给定假设中的给定公式）的推导。

一个更好的方法是使用你手工寻找推导时会用到的朴素方法：你从目标相继式开始，选中一个公式，然后“逆向”应用一条推理规则，生成一个或多个前提；如果这些前提有推导，它们就能与最后这条推理步骤组合成给定相继式的一个推导。与你在自然演绎系统中寻找推导的方法类似（尽管更复杂）。

但这个方法并非万无一失：如果你选错了规则，或者选错了公式的顺序，你可能会走进死胡同。它也不是完全确定的：在每一步，可能有不只一个公式可供选择，或者不只一条规则可以逆向应用。一个\emph{证明搜索算法}是一个完全确定的过程，它能在推导存在时（即万无一失地）找到一个推导。

有些推导系统比其他系统更适合简单的证明搜索算法。在相继式系统中，一个公式的主联结词，以及它出现在左边还是右边，决定了应该逆向应用哪条规则。例如，如果你想寻找 $\Gamma \Sequent \Delta, !A \land !B$ 的一个推导，其中 $!A \land !B$ 出现在右边，你应该逆向应用 \RightR{\land}，并生成两个相应的前提 $\Gamma \Sequent \Delta, !A$ 和 $\Gamma \Sequent \Delta, !B$。然而，如果你的系统是 \Log{G1c}，收缩规则的存在会使事情变得更困难，因为它给了你更多选择：你也可以逆向应用 \RightR{\Contraction} 并生成前提 $\Gamma \Sequent \Delta, !A \land
!B, !A \land !B$ 作为另一种选择。然后还可以生成前提 $\Gamma \Sequent
\Delta, !A \land !B, !A \land !B, !A \land !B$，如此等等。在 \Log{G2c} 中情况甚至更糟。逆向应用规则 \RightR{\land} 还意味着你必须猜测 $\Gamma_1$, $\Gamma_2$ 使得 $\Gamma = \Gamma_1 \cup \Gamma_2$，以及 $\Delta_1$, $\Delta_2$ 使得 $\Delta = \Delta_1 \cup \Delta_2$，然后去尝试前提 $\Gamma_1 \Sequent \Delta_1, !A$ 和 $\Gamma_2 \Sequent \Delta_2, !B$。一个错误的猜测可能导致你后来走入死胡同，于是你必须回溯到起点，选择不同的 $\Gamma_1$, $\Gamma_2$, $\Delta_1$, $\Delta_2$。如果公式是 $!A \lor !B$，你必须为 \RightR{\lor} 从两种可能性中选择一种，产生两个可能前提中的一个：$\Gamma \Sequent \Delta, !A$ 或 $\Gamma \Sequent \Delta, !B$。错误的选择同样将导致必须回溯。

系统 \Log{G3c} 没有这些问题。它没有结构规则，并且推理规则的选择完全由主联结词和公式出现的侧决定。因此，\Log{G3c} 比 \Log{G1c} 或 \Log{G2c} 更适合证明搜索。一次成功的搜索将产生一个 \Log{G3c} 推导，而这个推导随后可以被翻译成 \Log{G1c} 或 \Log{G2c}（就此而言，也可以是 \Log{N1c}）。所以，只要有了针对 \Log{G3c} 的证明搜索算法，以及将 \Log{G3c} 推导转换为其他系统推导的算法，也就解决了这些系统的证明搜索问题。

我们怎么知道一个证明搜索算法是万无一失的呢？这需要证明：如果算法没能找到推导，那么根本就不存在推导。毕竟，如果我们选了一个糟糕的算法，就有可能遇到这种情况：即使证明存在，算法也找不到。这种简单粗暴的算法不会有这个问题，因为它会遍历所有可能的推导。但一个证明搜索算法应该是高效的，并且付出最少的努力。

证明一个证明搜索算法万无一失的关键思路在于表明：如果算法失败了，我们可以从其失败的方式中推断出该相继式不可能有推导。而证明这点的方法就是证明该相继式不是有效的。在经典一阶逻辑中，有效的相继式 $\Gamma \Sequent \Delta$ 是指对任何结构，$\Gamma$ 中至少有一个公式为假，或者 $\Delta$ 中至少有一个公式为真。\Log{G3c} 是\emph{可靠的}，也就是说，它只能推导出有效的相继式。这是通过对推导的结构作归纳来建立的：公理显然是有效的，并且如果一个规则的前提是有效的，那么其结论也是有效的。为了证明一个证明搜索算法是万无一失的，我们要证明：如果对 $\Gamma \Sequent \Delta$ 的推导搜索失败了，我们可以定义一个结构，使得 $\Gamma$ 中的所有公式都为真，而 $\Delta$ 中的所有公式都为假。这同样也确立了 \Log{G3c} 是\emph{完备的}，即如果 $\Gamma \Sequent \Delta$ 是有效的，那么就存在一个它的 \Log{G3c} 推导。由于所有失败的证明搜索都表明 $\Gamma \Sequent \Delta$ 是无效的，所以对任何有效相继式的证明搜索都必定成功。

\end{document}